\chapter{Inner Products and Projections}\label{linear:ch:inner}

\begin{center}\itshape
Of all the principles that can be proposed for this purpose, I think there is none more general, more exact, or easier
to apply than the one we have used in this work: it consists of making the sum of the squares of the errors a minimum.
\upshape(Adrien-Marie Legendre, \emph{Nouvelles m\'ethodes pour la d\'etermination des orbites des com\`etes}, 1805)
\end{center}
\noindent\emph{Use Volume~1 and Appendices~\ref{linear:ch:duality} and~\ref{linear:ch:volume}. These finite-dimensional arguments complement the main Hilbert, spectral and singular-value chapters and supply tools for Appendix~\ref{linear:ch:inequalities}.}

Piazzi discovered Ceres on 1 January 1801 and followed it for about six weeks before sunlight hid it. A short noisy arc was all astronomers had for predicting where the new object would reappear. Gauss's prediction differed markedly from the others, and von Zach recovered Ceres near the predicted position at the end of that year.

Least squares became part of the ensuing history. Legendre published the principle in 1805: choose parameters making the sum of squared errors smallest. Gauss published his orbital methods in 1809 and claimed he had used least squares since 1795, provoking a priority dispute. The practical problem gives a clear geometric question: which vector in a model subspace is closest to the observations?

An inner product makes that closest vector an orthogonal projection. Repeated projection gives Gram--Schmidt and $QR$; its residual lengths give volumes and Gram determinants. Applied to monomials, it produces orthogonal polynomials, whose zeros are the nodes of Gauss quadrature. The same geometry links error fitting, polynomial approximation and integration.

\begin{think}
Fit a line $y=a+bx$ to $(0,0),(1,1),(2,1)$ by minimising squared vertical errors. Does it pass above or below the middle point? Should the three errors sum to zero? Example~\ref{linear:exa:line} computes the fit and explains the cancellation.
\end{think}

\section{Inner products and the Cauchy--Schwarz inequality}

Work over $\mathbb K=\R$ or $\C$, using conjugation also over $\R$, where it does nothing.
\begin{definition}\label{linear:def:inner}
An \emph{inner product} is a map $\langle\cdot,\cdot\rangle:V\times V\to\mathbb K$ linear in its first argument, conjugate symmetric, and positive definite: $\langle x,x\rangle>0$ for nonzero $x$. Its norm is $\|x\|=\sqrt{\langle x,x\rangle}$. Vectors with pairing zero are \emph{orthogonal}, written $x\perp y$.
\end{definition}

Conjugate symmetry makes the second argument conjugate-linear. On $\mathbb K^n$ use $\langle x,y\rangle=\sum_ix_i\bar y_i=y^*x$, where $y^*=\bar y^{\mathsf T}$. On continuous functions over a nondegenerate interval $[a,b]$, $a<b$, a positive continuous weight gives $\langle f,g\rangle=\int_a^bf\bar g\,w$. A continuous nonzero function is bounded away from zero on some subinterval, so its squared norm is positive. Restriction gives the polynomial example as well.

Expanding the norm gives $\|x+y\|^2=\|x\|^2+2\Re\langle x,y\rangle+\|y\|^2$. Orthogonality removes the middle term, giving Pythagoras.

\begin{theorem}[Cauchy--Schwarz]\label{linear:thm:cs}
For any $x,y$, $|\langle x,y\rangle|\le\|x\|\|y\|$, with equality exactly when the vectors are linearly dependent.
\end{theorem}
\begin{proof}
If $y=0$ the assertion and its equality condition hold. Otherwise put $p=\langle x,y\rangle y/\|y\|^2$. Direct pairing gives $x-p\perp y$, hence $x-p\perp p$. Pythagoras yields $\|x\|^2=\|p\|^2+\|x-p\|^2\ge|\langle x,y\rangle|^2/\|y\|^2$. Equality requires $x=p$, precisely dependence with nonzero $y$.
\end{proof}

\begin{corollary}\label{linear:cor:triangle}
The triangle inequality $\|x+y\|\le\|x\|+\|y\|$ holds, so $d(x,y)=\|x-y\|$ is a metric. For nonzero real-space vectors, the angle $\theta\in[0,\pi]$ is well defined by $\cos\theta=\langle x,y\rangle/(\|x\|\|y\|)$.
\end{corollary}
\begin{proof}
The norm expansion and Cauchy--Schwarz bound its square by $(\|x\|+\|y\|)^2$. Positivity, symmetry and homogeneity give the other metric properties. In the real case Cauchy--Schwarz puts the displayed quotient in $[-1,1]$, where cosine has a unique inverse in $[0,\pi]$.
\end{proof}

The real coordinate and integral inequalities $|\sum x_iy_i|^2\le\sum x_i^2\sum y_i^2$ and $|\int fg|^2\le\int f^2\int g^2$ are instances of the same statement. Complex versions use conjugates and absolute squares.

\begin{remark}
Jordan--von Neumann's criterion says a norm comes from an inner product precisely when it satisfies $\|x+y\|^2+\|x-y\|^2=2\|x\|^2+2\|y\|^2$. Expansion proves necessity; Exercise~\ref{linear:ex:jvn} proves sufficiency over $\R$. The norm $|x_1|+|x_2|$ fails this parallelogram law at the two coordinate unit vectors.
\end{remark}

\section{Orthonormal bases and Gram--Schmidt}

An \emph{orthonormal family} satisfies $\langle e_i,e_j\rangle=\delta_{ij}$.
\begin{proposition}\label{linear:prop:orthfamily}
Such a family is independent. In its span, $x=\sum_ix_ie_i$ has coordinates $x_i=\langle x,e_i\rangle$ and norm squared $\sum_i|x_i|^2$.
\end{proposition}
\begin{proof}
Pair the expansion with each $e_i$ to recover its coefficient. A zero vector therefore has only zero coefficients. Pairing the full expansion with itself cancels all off-diagonal terms, giving the norm formula.
\end{proof}

Orthogonal coordinates need only inner products, rather than a linear system. The next construction finds such coordinates while preserving successive spans.

\begin{theorem}[Gram--Schmidt]\label{linear:thm:gramschmidt}
For independent $v_1,\ldots,v_k$, there is exactly one orthonormal family $e_1,\ldots,e_k$ with matching successive spans and $\langle v_j,e_j\rangle>0$ for every $j$. It is given by
\[u_j=v_j-\sum_{i<j}\langle v_j,e_i\rangle e_i,\qquad e_j=u_j/\|u_j\|.\]
\end{theorem}
\begin{proof}
At each stage the subtraction makes $u_j$ orthogonal to earlier $e_i$. Independence makes it nonzero, and normalisation gives the desired span and positive pairing $\|u_j\|$. For uniqueness, the complement of the earlier span inside the new span is a line. Any other unit choice is $ce_j$ with $|c|=1$; its required pairing is $\bar c\|u_j\|>0$, forcing $c=1$.
\end{proof}

Consequently every finite-dimensional inner-product space has an orthonormal basis, and every orthonormal family in it extends to one. Complete it first to a basis and run the construction on the added vectors.

\begin{corollary}[$QR$ factorisation]\label{linear:cor:qr}
A rectangular matrix with independent columns has a unique factorisation $A=QR$, where $Q^*Q=I$ and $R$ is square upper triangular with positive diagonal.
\end{corollary}
\begin{proof}
Apply Gram--Schmidt to the columns. Each $a_j=\sum_{i\le j}\langle a_j,e_i\rangle e_i$, giving $Q$ and $r_{ij}$, with $r_{jj}=\|u_j\|>0$. Conversely any such factorisation has the same successive column spans and positive diagonal pairings. Gram--Schmidt uniqueness therefore fixes $Q$, and $R=Q^*A$ is fixed as well.
\end{proof}

\section{Orthogonal projections}

For any subset $S$, let $S^\perp$ be the vectors orthogonal to every element of $S$.
\begin{theorem}[orthogonal projection]\label{linear:thm:projection}
Let $W$ be a finite-dimensional subspace of an inner-product space $V$, which need not be finite-dimensional. For an orthonormal basis $e_1,\ldots,e_k$ of $W$, put $P_Wx=\sum_i\langle x,e_i\rangle e_i$. Then $x-P_Wx\in W^\perp$ and $V=W\oplus W^\perp$. Moreover
\[\|x-w\|^2=\|x-P_Wx\|^2+\|P_Wx-w\|^2\qquad(w\in W),\]
so $P_Wx$ is the unique nearest point. The map is linear and idempotent, and $\langle P_Wx,y\rangle=\langle x,P_Wy\rangle$.
\end{theorem}
\begin{proof}
Pairing the residual with each basis vector gives zero. A vector in $W\cap W^\perp$ has zero squared norm, proving uniqueness of the decomposition. The displayed identity is Pythagoras for its two orthogonal terms, so minimisation and uniqueness follow. The formula is linear, fixes $W$ and thus is idempotent. Finally both pairings in the last identity equal $\sum_i\langle x,e_i\rangle\overline{\langle y,e_i\rangle}$. Figure~\ref{linear:fig:projection} gives the geometric picture.
\end{proof}
\begin{figure}[ht]
\centering
\fig{linear-ch04-projection}
\caption{The orthogonal projection. Among all points $w$ of the plane $W$, the foot $P_Wx$ of the perpendicular is the
closest to $x$, because $x-w$ is the hypotenuse of a right triangle whose other side is $x-P_Wx$.}
\label{linear:fig:projection}
\end{figure}

Pythagoras for the projection gives Bessel's inequality:
\begin{equation}\label{linear:eq:bessel}
\sum_i|\langle x,e_i\rangle|^2\le\|x\|^2.
\end{equation}
For an infinite orthonormal family interpret this sum as the supremum of finite subsums. Each finite projection gives the bound. Main Chapter~\ref{ch:hilbert} develops completeness and Parseval equality; Exercise~\ref{linear:ex:basel} uses just the inequality.

In finite dimension, $\dim W^\perp=\dim V-\dim W$ and $W^{\perp\perp}=W$. Compare annihilators in Appendix~\ref{linear:ch:duality}. The inner product will identify vectors with forms, conjugate-linearly in the complex case, in Section~\ref{linear:sec:adjoint}.

\section{Least squares}

For rectangular $A$ and observed $b$, the equation $Ax=b$ may be inconsistent. Minimising $\|Ax-b\|$ finds the closest vector in the column space, regardless of whether the number of rows exceeds the number of columns.

\begin{theorem}[least squares]\label{linear:thm:lsq}
A vector $x$ minimises $\|Ax-b\|$ exactly when it satisfies $A^*Ax=A^*b$. Minimisers always exist and are unique exactly when $A$ has independent columns. In that case
\[x=(A^*A)^{-1}A^*b=R^{-1}Q^*b,\qquad A=QR.\]
\end{theorem}
\begin{proof}
The optimal image is $P_{\operatorname{Im}A}b$. Thus minimisation is equivalent to $b-Ax\perp\operatorname{Im}A$, or $y^*A^*(b-Ax)=0$ for all $y$, giving the normal equations. An optimal image has a preimage, proving existence. All its preimages differ by $\ker A$, proving the uniqueness criterion. With independent columns, $x^*A^*Ax=\|Ax\|^2$ proves $A^*A$ invertible. Substituting $A=QR$ into the equations and cancelling $R^*$ gives the second formula.
\end{proof}

\begin{example}\label{linear:exa:line}
For the opening data, the design matrix and observation vector are
\[A=\begin{pmatrix}1&0\\1&1\\1&2\end{pmatrix},\qquad y=\begin{pmatrix}0\\1\\1\end{pmatrix}.
\]
Their normal equations are $3a+3b=2$, $3a+5b=3$, so $a=1/6,b=1/2$. The residuals are $-1/6,1/3,-1/6$, summing to zero. The line lies below the middle point, as Figure~\ref{linear:fig:lsq} shows. The zero sum follows from residual orthogonality to the constant column; every least-squares model containing constants has this property.
\end{example}
\begin{figure}[ht]
\centering
\fig{linear-ch04-lsq}
\caption{The least squares line of Example~\ref{linear:exa:line}. The sum of the squares of the rust segments is $\frac16$,
smaller than for any other line.}
\label{linear:fig:lsq}
\end{figure}


Squared errors give linear normal equations. They also have a statistical optimality property without assuming Gaussian noise. Suppose real observations satisfy $b=Ax_0+\varepsilon$, with $\mathbb E\varepsilon=0$ and covariance $\sigma^2I$. A linear estimator $Cb$ is unbiased for every $x_0$ exactly when $CA=I$; its covariance is then $\sigma^2CC^{\mathsf T}$. Equal variances and zero pairwise covariances suffice; independence is stronger than needed.

\begin{theorem}[Gauss--Markov]\label{linear:thm:gaussmarkov}
For real $A$ with independent columns, let $C_0=(A^{\mathsf T}A)^{-1}A^{\mathsf T}$. Every $C$ with $CA=I$ satisfies $CC^{\mathsf T}-C_0C_0^{\mathsf T}\ge0$ as a positive semidefinite matrix. Thus every scalar linear estimate $u^{\mathsf T}Cb$ has variance at least that of $u^{\mathsf T}C_0b$.
\end{theorem}
\begin{proof}
Set $D=C-C_0$. Then $DA=0$, so $DC_0^{\mathsf T}=DA(A^{\mathsf T}A)^{-1}=0$ and its transpose also vanishes. Expanding gives $CC^{\mathsf T}=C_0C_0^{\mathsf T}+DD^{\mathsf T}$. Its difference has quadratic form $\|D^{\mathsf T}u\|^2\ge0$. Multiply by $\sigma^2$ for the variance comparison.
\end{proof}

The extra estimator component contributes only the non-negative term $DD^{\mathsf T}$. This is the same orthogonal-error geometry as the closest-point theorem.

\section{Gram determinants and volumes}

For vectors $v_1,\ldots,v_k$, define their Gram matrix by $G_{ij}=\langle v_j,v_i\rangle$ and their Gram determinant by $\det G$. With coordinate columns $A$ in an orthonormal basis, $G=A^*A$.

\begin{theorem}[Gram determinants]\label{linear:thm:gram}
Put $W_{j-1}=\operatorname{span}(v_1,\ldots,v_{j-1})$ and $u_j=v_j-P_{W_{j-1}}v_j$.
\begin{enumerate}[label=(\alph*)]
\item The Gram determinant equals $\prod_{j=1}^k\|u_j\|^2$. It is non-negative and vanishes exactly for dependent vectors.
\item For independent $v_1,\ldots,v_k$ spanning $W$,
\[d(x,W)^2=\frac{\operatorname{Gram}(v_1,\ldots,v_k,x)}{\operatorname{Gram}(v_1,\ldots,v_k)}.\]
\item For real columns $A$ in $\R^n$, $\operatorname{Gram}(v_1,\ldots,v_k)=\sum_{|I|=k}\det(A_I)^2$.
\end{enumerate}
\end{theorem}
\begin{proof}
(a) Dependence gives a nonzero coefficient vector killed by $A$ and therefore by $A^*A$, so its determinant is zero; the first dependent residual is also zero. For independent columns, $A=QR$ gives $G=R^*R$ and $\det G=|\det R|^2=\prod_j\|u_j\|^2$. (b) Append $x$ in (a); the last residual is $x-P_Wx$, while the earlier factors are unchanged. (c) Apply real Cauchy--Binet, Corollary~\ref{linear:cor:pythagoras}.
\end{proof}

In real geometry, successive perpendicular residual lengths are successive base and height factors. Their product is the parallelepiped's volume, so the Gram determinant is its square. The two-vector case is Cauchy--Schwarz again; the appended-vector ratio computes distance by determinants.

\section{The Cauchy determinant and the Hilbert matrix}
\begin{theorem}[Cauchy, 1841]\label{linear:thm:cauchydet}
Whenever all $x_i+y_j$ are nonzero,
\[\det\left(\frac1{x_i+y_j}\right)_{i,j=1}^n
 =\frac{\prod_{i<j}(x_j-x_i)(y_j-y_i)}{\prod_{i,j}(x_i+y_j)}.\]
\end{theorem}
\begin{proof}
If two $x$ values or two $y$ values coincide, both sides vanish. Otherwise induct on $n$. Subtract row $n$ from each earlier row, giving differences $(x_n-x_i)/((x_i+y_j)(x_n+y_j))$. Factor $x_n-x_i$ from each such row and $(x_n+y_j)^{-1}$ from every column. The last row becomes all ones and the earlier rows return to Cauchy form.

Subtract the last column from the others. The earlier entries become $(y_n-y_j)/((x_i+y_j)(x_i+y_n))$, and the last row becomes $(0,\ldots,0,1)$. Factor those numerators from columns and $(x_i+y_n)^{-1}$ from earlier rows, then expand along the last row. This gives
\[\det C=\frac{\prod_{i<n}(x_n-x_i)\prod_{j<n}(y_n-y_j)}{\prod_{j\le n}(x_n+y_j)\prod_{i<n}(x_i+y_n)}\det C_{n-1}.
\]
These are exactly the new factors in the claimed product, proving the induction from $n=1$.
\end{proof}

Hilbert's 1894 approximation problem led to the Gram matrix of $1,t,\ldots,t^{n-1}$ on $[0,1]$:
\[H_n=\left(\int_0^1t^{i+j}\dd t\right)_{i,j=0}^{n-1}
 =\left(\frac1{i+j+1}\right)_{i,j=0}^{n-1}.
\]
Independent monomials make it positive definite, despite the difficulty of seeing that directly from its entries.

\begin{corollary}\label{linear:cor:hilbert}
For $c_n=\prod_{k=1}^{n-1}k!$, one has $\det H_n=c_n^4/c_{2n}$. The $L^2[0,1]$ distance from $t^n$ to lower-degree polynomials is
\[d(t^n,\R[t]_{<n})=\frac{(n!)^2}{(2n)!\sqrt{2n+1}}
 =\frac1{\binom{2n}n\sqrt{2n+1}}.\]
\end{corollary}
\begin{proof}
Apply Cauchy's formula with $x_i=i,y_j=j+1$. Adding the next row and column gives
\[\frac{\det H_{n+1}}{\det H_n}=\frac{(n!)^4}{(2n)!(2n+1)!}.
\]
The quotient of consecutive $c_n^4/c_{2n}$ has the same value and starts at $\det H_1=1$. The Gram ratio theorem gives squared distance equal to this determinant ratio, whose square root is the displayed expression. At $n=0$, the distance from one to the zero polynomial space is one, also agreeing with the distance formula.
\end{proof}

For degrees one and two, the distances are $1/(2\sqrt3)$ and $1/(6\sqrt5)$. Stirling applied to the exact factorial product shows $\log\det H_n=-n^2\log4+O(n\log n)$. The monomials approach dependence rapidly, and Hilbert matrices are correspondingly demanding numerical examples. Their reciprocal determinants are nevertheless integers, as Exercise~\ref{linear:ex:hilbertint} proves.

\begin{remark}
Hilbert used lattice-point arguments to find nonzero integer polynomials with arbitrarily small $L^2$ integral on intervals shorter than four. That application needs Minkowski's theorem, proved later in Volume~13, Chapter~5; we quote it here and use no lattice-point result below. Borwein--Erd\'elyi discuss weighted and uniform-norm versions of the integer Chebyshev problem.
\end{remark}

\section{Orthogonal polynomials}

Let $I$ have nonempty interior and let $w$ be positive and continuous there, with $\int_I|t|^kw(t)\dd t<\infty$ for every non-negative integer $k$. Then $\langle f,g\rangle=\int_Ifgw$ is an inner product on real polynomials: a nonzero polynomial cannot vanish throughout an interval. Gram--Schmidt on the monomials, rescaled to leading coefficient one, gives unique monic polynomials $p_n$ of degree $n$ orthogonal to all lower degrees, with $p_0=1$.

\begin{example}[Legendre polynomials]\label{linear:exa:legendre}
For $I=[-1,1]$ and $w=1$, Rodrigues's 1816 formula is
\[P_n(t)=\frac1{2^nn!}\frac{\dd^n}{\dd t^n}(t^2-1)^n.
\]
Every derivative of $(t^2-1)^n$ below order $n$ vanishes at both endpoints. Integrating $\int t^kP_n$ by parts $n$ times therefore gives zero when $k<n$, proving orthogonality. Exercise~\ref{linear:ex:legendrenorm} verifies $P_n(1)=1$, leading coefficient $\binom{2n}n/2^n$ and norm squared $2/(2n+1)$. Division by the leading coefficient gives the monic version. Other choices include Chebyshev's weight $(1-t^2)^{-1/2}$ on $(-1,1)$, Laguerre's $e^{-t}$ on $[0,\infty[$ and Hermite's $e^{-t^2}$ on $\R$.
\end{example}
\begin{figure}[ht]
\centering
\fig{linear-ch04-legendre}
\caption{The Legendre polynomials $P_0,\dots,P_5$. Each $P_n$ has $n$ simple zeros in $(-1,1)$, and the zeros of
consecutive polynomials interlace.}
\label{linear:fig:legendre}
\end{figure}

Multiplication by the variable couples only neighbouring orthogonal degrees.
\begin{theorem}[three-term recurrence]\label{linear:thm:recurrence}
There are real $\alpha_n$ and positive $\beta_n$ such that, with $p_{-1}=0$,
\[p_{n+1}(t)=(t-\alpha_n)p_n(t)-\beta_np_{n-1}(t),\]
where $\alpha_n=\langle tp_n,p_n\rangle/\|p_n\|^2$ and $\beta_n=\|p_n\|^2/\|p_{n-1}\|^2$ for $n\ge1$. At $n=0$, choose any $\beta_0>0$, since it multiplies zero.
\end{theorem}
\begin{proof}
The polynomial $tp_n-p_{n+1}$ has degree at most $n$. Its coefficient along $p_k$ is $\langle tp_n,p_k\rangle/\|p_k\|^2$. Symmetry of multiplication makes this $\langle p_n,tp_k\rangle/\|p_k\|^2$, zero when $k<n-1$. The remaining diagonal coefficient is $\alpha_n$. At $k=n-1$, $tp_{n-1}-p_n$ has lower degree, so the numerator is $\|p_n\|^2$, giving $\beta_n$. Its positivity follows from the positive norms.
\end{proof}

\begin{theorem}[zeros]\label{linear:thm:zeros}
The degree-$n$ polynomial $p_n$ has $n$ distinct real roots, all in the interior of $I$.
\end{theorem}
\begin{proof}
Let $t_1,\ldots,t_k$ be its sign-changing interior roots and set $q(t)=\prod_i(t-t_i)$, or one if there are none. The product $p_nq$ has one fixed sign on $I$, is nonzero on a subinterval and is integrable against the positive weight. Hence $\langle p_n,q\rangle\ne0$. Orthogonality would force zero if $k<n$. Thus $k\ge n$, which uses all available degree and makes every root simple and interior.
\end{proof}

\begin{proposition}[the Jacobi matrix]\label{linear:prop:jacobi}
The roots of $p_n$ are exactly the eigenvalues of the symmetric tridiagonal matrix with diagonal $\alpha_0,\ldots,\alpha_{n-1}$ and adjacent entries $\sqrt{\beta_1},\ldots,\sqrt{\beta_{n-1}}$.
\end{proposition}
\begin{proof}
Let $T$ instead have superdiagonal ones and subdiagonal $\beta_k$. The recurrence gives $Tv(t)=tv(t)-p_n(t)e_n$, where $v(t)=(p_0(t),\ldots,p_{n-1}(t))^{\mathsf T}$. Every root gives a nonzero eigenvector because $p_0=1$. There are $n$ distinct roots, so $\chi_T=p_n$. With $D=\operatorname{diag}(1,\sqrt{\beta_1},\sqrt{\beta_1\beta_2},\ldots)$, direct multiplication gives $T=DJ_nD^{-1}$ for the stated symmetric matrix. Thus their characteristic polynomials agree.
\end{proof}

Nested Jacobi matrices give weak root interlacing by Chapter~\ref{linear:ch:spectral}'s Cauchy theorem. It is strict: a shared root of $p_n,p_{n+1}$ would, by the recurrence and $\beta_n>0$, also kill $p_{n-1}$ and then every preceding polynomial, contradicting $p_0=1$.

\section{Gauss quadrature}

Fixed distinct nodes give an interpolatory rule exact below degree $n$, with weights $\lambda_i=\int_I\ell_iw$. Choosing the nodes as orthogonal-polynomial zeros raises exactness to degree $2n-1$.

\begin{theorem}[Gauss quadrature]\label{linear:thm:gaussquad}
For $n\ge1$, take the roots $t_i$ of $p_n$ and $\lambda_i=\int_I\ell_iw$. Then
\[\int_Ifw=\sum_{i=1}^n\lambda_if(t_i)\qquad(\deg f\le2n-1).
\]
Each weight is $\lambda_i=\int_I\ell_i^2w>0$. No rule with $n$ real nodes is exact for all degree-$2n$ polynomials.
\end{theorem}
\begin{proof}
Divide $f=qp_n+r$, with both $q,r$ of degree at most $n-1$. Orthogonality kills the integral of $qp_n$, and its node values vanish too. Interpolation integrates $r$ exactly, proving the identity. Apply it to $\ell_i^2$ to obtain the positive-weight formula. For any proposed $n$ real nodes $s_i$, the polynomial $\prod_i(t-s_i)^2$ has positive integral but zero node values, proving the degree obstruction.
\end{proof}

For two-node Legendre quadrature, the roots are $\pm1/\sqrt3$ and both weights are one. Two evaluations thus integrate any cubic on $[-1,1]$ exactly. Positive weights also control uniform approximation errors and give convergence for functions continuous on a compact interval (Exercise~\ref{linear:ex:gaussconv}).

\section{Adjoints and isometries}\label{linear:sec:adjoint}

Assume finite dimension for the rest of this section.
\begin{theorem}[Riesz representation]\label{linear:thm:riesz}
Every linear form $\varphi$ is represented uniquely by a vector $y$ with $\varphi(x)=\langle x,y\rangle$ for all $x$.
\end{theorem}
\begin{proof}
For an orthonormal basis set $y=\sum_i\overline{\varphi(e_i)}e_i$. Expansion verifies the identity. Two representatives have difference orthogonal to every vector, including itself, and so agree.
\end{proof}

The resulting identification with the dual is conjugate-linear over $\C$ and linear over $\R$. For an operator $A$, apply it to the form $x\mapsto\langle Ax,y\rangle$.
\begin{proposition}[the adjoint]\label{linear:prop:adjoint}
There is exactly one linear $A^*$ with $\langle Ax,y\rangle=\langle x,A^*y\rangle$. Its orthonormal-coordinate matrix is conjugate transpose. It satisfies $(AB)^*=B^*A^*$, $A^{**}=A$, $\ker A^*=(\operatorname{Im}A)^\perp$ and $\operatorname{Im}A^*=(\ker A)^\perp$.
\end{proposition}
\begin{proof}
Riesz gives the unique representing vector for each $y$. Conjugate-linearity in both second arguments makes their assignments linear in $y$. Evaluating at basis vectors gives entry $\overline{a_{ji}}$. Applying the defining equality twice proves the product and double-adjoint rules. Finally $A^*y=0$ exactly when $\langle Ax,y\rangle=0$ for all $x$, proving the kernel equality. Apply it to $A^*$ and take orthogonal complements for the image equality.
\end{proof}

This matches the transpose construction after identifying the dual with the original space. Fredholm's condition becomes orthogonality of $b$ to $\ker A^*$.

\begin{theorem}[unitary operators]\label{linear:thm:unitary}
For a linear operator $U$ in finite dimension, the following are equivalent: preserving inner products; preserving norms; carrying some orthonormal basis, hence every one, to an orthonormal basis; and $U^*U=I$.
\end{theorem}
\begin{proof}
Inner-product preservation gives norm preservation. Conversely polarisation recovers the real part by $(\|x+y\|^2-\|x-y\|^2)/4$ and, over $\C$, the imaginary part as $\Re\langle x,iy\rangle$. Thus norms determine all pairings for a linear map. Preserving one orthonormal basis preserves pairings by coordinate expansion, and preserving all pairings clearly preserves every such basis. The adjoint identity makes inner-product preservation equivalent to $U^*U=I$.
\end{proof}

These operators are unitary over $\C$ and orthogonal over $\R$. Their matrices have orthonormal columns and form a group; their inverse is the adjoint. In real Euclidean space, metric preservation already forces linearity once the origin is fixed.

\begin{theorem}\label{linear:thm:isometry}
Any map $f:\R^n\to\R^n$ fixing zero and preserving all distances is linear and orthogonal.
\end{theorem}
\begin{proof}
Distance preservation from zero gives norms. The identity $\langle x,y\rangle=(\|x\|^2+\|y\|^2-\|x-y\|^2)/2$ therefore gives pairing preservation. Expand the squared norm of $f(x+y)-f(x)-f(y)$ into these pairings; it equals $\|(x+y)-x-y\|^2=0$. Likewise $\|f(\lambda x)-\lambda f(x)\|^2=0$ for real $\lambda$. Thus $f$ is real-linear and preserves inner products.
\end{proof}

Every Euclidean rigid motion is consequently a translation followed by an orthogonal map. Conversely, positive definite Hermitian matrices describe all inner products on coordinate space through $\langle x,y\rangle_G=y^*Gx$.

\begin{corollary}[Cholesky factorisation]\label{linear:cor:cholesky}
Every positive definite matrix has a unique factorisation $G=R^*R$ with $R$ upper triangular and positive on its diagonal.
\end{corollary}
\begin{proof}
Run Gram--Schmidt on the coordinate basis for the $G$ inner product. If its new basis matrix is $F$ and its coefficient matrix $R$, then $I=FR$ and $F^*GF=I$. Hence $G=R^*R$. If also $G=R'^*R'$, the matrix $R'R^{-1}$ is unitary and upper triangular with positive diagonal. Its adjoint, also its inverse, is both lower and upper triangular, so it is diagonal. Its positive unit-modulus entries are all one, proving uniqueness.
\end{proof}

Appendix~\ref{linear:ch:inequalities} uses this factorisation for positive-matrix inequalities.

\begin{beauty}
Repeatedly remove a perpendicular component from the monomials. Their residual norms determine the Hilbert determinants, while the resulting orthogonal polynomials give tridiagonal eigenvalue problems and optimal quadrature nodes. The same projection calculation organised the original least-squares fit. These are different uses of one precise orthogonality identity.
\end{beauty}

\begin{figure}[ht]
\centering\fig{appD-rw-weights}
\caption{Weighted least squares at the three opening abscissae, with weights $1,2,1$, gives residuals $-1/4,1/4,-1/4$. Their weighted sum vanishes, although their ordinary sum does not.}
\label{fig:v8appDweights}
\end{figure}
\begin{figure}[ht]
\centering\fig{appD-rw-gauss}
\caption{The three-node Legendre rule uses nodes $-\sqrt{3/5},0,\sqrt{3/5}$ with weights $5/9,8/9,5/9$. It integrates every polynomial through degree five exactly; the sixth-degree error is $8/175$.}
\label{fig:v8appDgauss}
\end{figure}
\section{Exercises}
\begin{exercise}[$\star$]\label{linear:ex:polar}
(a) Prove the parallelogram norm identity and interpret its sides and diagonals geometrically. (b) For a complex inner product linear in the first argument, establish $\langle x,y\rangle=\tfrac14\sum_{k=0}^3i^k\|x+i^ky\|^2$.
\end{exercise}
\begin{exercise}[$\star\star\star$]\label{linear:ex:jvn}
For a real norm satisfying the parallelogram law, prove that $(\|x+y\|^2-\|x-y\|^2)/4$ defines an inner product inducing that norm.
\end{exercise}
\begin{exercise}[$\star$]\label{linear:ex:qr}
Compute $QR$ for $\begin{pmatrix}1&1&0\\1&0&1\\0&1&1\end{pmatrix}$ and verify the determinant magnitude equals the product of the positive $R$ diagonals.
\end{exercise}
\begin{exercise}[$\star$]\label{linear:ex:projmatrix}
For real $A$ with independent columns, prove $A(A^{\mathsf T}A)^{-1}A^{\mathsf T}$ projects orthogonally onto its column space. Conversely prove $P^2=P=P^{\mathsf T}$ characterises projection onto $\operatorname{Im}P$, with $\operatorname{tr}P=\operatorname{rank}P$.
\end{exercise}
\begin{exercise}[$\star\star$]\label{linear:ex:parabola}
Fit a quadratic by least squares to $(-2,3),(-1,1),(0,0),(1,1),(2,4)$. Repeat in the basis $1,x,x^2-2$ and explain its diagonal normal equations.
\end{exercise}
\begin{exercise}[$\star$]\label{linear:ex:distpoly}
Find the affine polynomial nearest $t^2$ in $L^2[0,1]$. Compute its squared error and compare with Corollary~\ref{linear:cor:hilbert}.
\end{exercise}
\begin{exercise}[$\star$]\label{linear:ex:hilbertint}
Rewrite the consecutive Hilbert determinant ratio as $((2n+1)\binom{2n}n^2)^{-1}$. Deduce integer reciprocal determinants and compute them through size four.
\end{exercise}
\begin{exercise}[$\star\star$]\label{linear:ex:obtuse}
Construct $n+1$ vectors in $\R^n$ with negative pairwise inner products, and exclude $n+2$. For the upper bound project away one vector and use dimension induction.
\end{exercise}
\begin{exercise}[$\star\star$]\label{linear:ex:basel}
In continuous real functions on $[-\pi,\pi]$ with integral inner product, normalise the sine functions to an orthonormal family. Apply Bessel to $f(t)=t$ to obtain $\sum_{k\ge1}k^{-2}\le\pi^2/6$. Parseval, Theorem~\ref{thm:fourierl2}, gives equality, but is unnecessary for this bound.
\end{exercise}
\begin{exercise}[$\star\star$]\label{linear:ex:legendrenorm}
From Rodrigues's formula prove the endpoint value one, leading coefficient $\binom{2n}n/2^n$, and squared norm $2/(2n+1)$. Establish the auxiliary integral $\int_{-1}^1(1-t^2)^n\dd t=2^{2n+1}(n!)^2/(2n+1)!$ rather than assuming it.
\end{exercise}
\begin{exercise}[$\star\star$]\label{linear:ex:chebyshev}
Define $T_n(\cos\theta)=\cos(n\theta)$. Prove polynomial degree $n$ and orthogonality for $(1-t^2)^{-1/2}$ on $(-1,1)$. Find its $n$-point Gauss rule: nodes $\cos((2k-1)\pi/(2n))$ and weights $\pi/n$.
\end{exercise}
\begin{exercise}[$\star\star$]\label{linear:ex:gaussconv}
For the Gauss rules $G_n$ of the given positive weight on a compact interval $I=[a,b]$, prove $G_n(f)\to\int_Ifw$ for every $f\in C([a,b])$ by positive weights and Weierstrass. Continuity only on a bounded open interval does not imply uniform polynomial approximation; the endpoint condition here makes the intended hypothesis explicit.
\end{exercise}
\begin{exercise}[$\star\star$]\label{linear:ex:householder}
For nonzero real $u$, prove $H_u=I-2uu^{\mathsf T}/(u^{\mathsf T}u)$ is a symmetric orthogonal reflection. Show $H_{x-y}x=y$ for distinct equal-norm $x,y$. Then factor every orthogonal $n\times n$ matrix into at most $n$ such reflections.
\end{exercise}
\begin{exercise}[$\star\star\star$]\label{linear:ex:muntz}
Let $0\le a_1<a_2<\cdots$ and let integer $m\ge0$ differ from every $a_k$. Prove that the distance from $t^m$ to the first $N$ powers in $L^2[0,1]$ satisfies
\[d_N^2=\frac1{2m+1}\prod_{k=1}^N\left(\frac{a_k-m}{a_k+m+1}\right)^2.\]
If $a_k\to\infty$ and the reciprocal sum over the positive $a_k$ diverges, deduce $d_N\to0$. An initial zero exponent is omitted from the reciprocal sum; finite initial terms do not affect its divergence.
\end{exercise}

\begin{exercise}[$\star\star$]\label{ex:v8cDa}
\emph{Changing observation weights.} Fit the opening three points using squared-error weights $1,2,1$. Derive the weighted normal equations by multiplying rows by square roots of the weights. Find the line, residuals and minimum error, and check their weighted zero-sum identity.
\end{exercise}
\begin{exercise}[$\star\star$]\label{ex:v8cDb}
\emph{A least-norm interpolant.} Among real polynomials of degree at most two satisfying $p(0)=p(1)=1$, minimise $\int_0^1p(t)^2\dd t$. Prove uniqueness and relate the minimising condition to orthogonality along the space of polynomials vanishing at both endpoints.
\end{exercise}
\begin{exercise}[$\star\star$]\label{ex:v8cDc}
\emph{Three Gauss nodes and the first error.} Derive the Legendre three-node weights $5/9,8/9,5/9$ at $-\sqrt{3/5},0,\sqrt{3/5}$. Prove exactness through degree five and compute the quadrature error for $t^6$ as $8/175$ using the monic orthogonal cubic.
\end{exercise}
\begin{exercise}[$\star\star$]\label{ex:v8cDd}
\emph{When two projections combine.} For orthogonal projections $P,Q$ in finite dimension, prove $PQ$ is an orthogonal projection exactly when $PQ=QP$. In that case identify its image as the intersection of the two images. Give two line projections in $\R^2$ whose product is not a projection.
\end{exercise}

\section{Hints}
\begin{hint}{linear:ex:polar}
Expand squared norms and sum the powers of $i$.
\end{hint}
\begin{hint}{linear:ex:jvn}
First prove $B(x+z,y)+B(x-z,y)=2B(x,y)$ for the proposed pairing. Doubling then gives additivity, and norm continuity extends rational homogeneity to real homogeneity.
\end{hint}
\begin{hint}{linear:ex:qr}
Subtract earlier orthogonal components column by column and retain positive residual lengths.
\end{hint}
\begin{hint}{linear:ex:projmatrix}
Use the least-squares image for the forward statement and residual orthogonality for the converse.
\end{hint}
\begin{hint}{linear:ex:parabola}
The five abscissae give zero sums for $x,x^3,x^2-2$ and $x(x^2-2)$.
\end{hint}
\begin{hint}{linear:ex:distpoly}
Make $t^2-a-bt$ orthogonal to both one and $t$.
\end{hint}
\begin{hint}{linear:ex:hilbertint}
Multiply the reciprocal ratios starting from determinant one.
\end{hint}
\begin{hint}{linear:ex:obtuse}
Projection onto the last vector's orthogonal complement subtracts a positive quantity from the remaining pairwise inner products.
\end{hint}
\begin{hint}{linear:ex:basel}
Integration by parts gives $\int_{-\pi}^\pi t\sin(kt)\dd t=2\pi(-1)^{k+1}/k$.
\end{hint}
\begin{hint}{linear:ex:legendrenorm}
At $t=1$ only one Leibniz term survives. For the norm replace one polynomial factor by its leading monomial, then integrate by parts $n$ times.
\end{hint}
\begin{hint}{linear:ex:chebyshev}
The cosine recurrence gives polynomiality. Substitution $t=\cos\theta$ gives orthogonality and a finite cosine sum for the weights.
\end{hint}
\begin{hint}{linear:ex:gaussconv}
Both the positive rule and the integral have norm bounded by the total weight mass. Apply them to a uniformly small polynomial-approximation error.
\end{hint}
\begin{hint}{linear:ex:householder}
First fix $e_1$ with one reflection, then induct on its orthogonal complement.
\end{hint}
\begin{hint}{linear:ex:muntz}
The power-function Gram entries are $1/(a+b+1)$. Add exponent $m$ to Cauchy's determinant, then estimate the logarithm of its product tail.
\end{hint}
\begin{hint}{ex:v8cDa}
Use $A^{\mathsf T}WA$ and $A^{\mathsf T}Wy$, with $W=\operatorname{diag}(1,2,1)$.
\end{hint}
\begin{hint}{ex:v8cDb}
Every feasible polynomial is $1+ct(1-t)$.
\end{hint}
\begin{hint}{ex:v8cDc}
Symmetry leaves two weights. Apply the rule to one and $t^2$, then compare $t^6$ with $(t^3-3t/5)^2$.
\end{hint}
\begin{hint}{ex:v8cDd}
The adjoint of $PQ$ is $QP$. For the counterexample use the coordinate axis and the line $y=x$.
\end{hint}

\section{Solutions}
\begin{solution}{linear:ex:polar}
\idea Norm expansions retain or cancel the inner-product cross terms.
\writeup (a) Add $\|x+y\|^2$ and $\|x-y\|^2$; the cross terms cancel. A parallelogram's diagonal vectors are $x+y,x-y$, while its four sides have squared lengths summing to $2\|x\|^2+2\|y\|^2$. (b) Expansion gives $\|x+i^ky\|^2=\|x\|^2+\|y\|^2+i^{-k}\langle x,y\rangle+i^k\langle y,x\rangle$. Multiplication by $i^k$ and summation cancel the constant-norm and $i^{2k}$ terms, leaving $4\langle x,y\rangle$.
\end{solution}
\begin{solution}{linear:ex:jvn}
\idea The parallelogram law forces the proposed polarisation form to be additive; continuity supplies real homogeneity.
\writeup Denote the proposed form by $B$. It is symmetric, $B(x,x)=\|x\|^2$ and $B(0,y)=0$. Apply the parallelogram law to $x+y,z$ and $x-y,z$ and subtract. Dividing by four gives $B(x+z,y)+B(x-z,y)=2B(x,y)$. Setting $z=x$ gives $B(2x,y)=2B(x,y)$. Now put $x=(u+v)/2,z=(u-v)/2$ to obtain additivity. Integer and then rational homogeneity follow. The formula for $B$ is continuous in each argument because the norm is continuous, so rational approximation proves real homogeneity. Symmetry gives bilinearity, and positive definiteness follows from $B(x,x)=\|x\|^2$.
\end{solution}
\begin{solution}{linear:ex:qr}
\idea Orthogonalise columns and store each removed component in $R$.
\writeup The unit columns are $e_1=(1,1,0)/\sqrt2$, $e_2=(1,-1,2)/\sqrt6$ and $e_3=(-1,1,1)/\sqrt3$. Their positive residual lengths are $\sqrt2,\sqrt{3/2},2/\sqrt3$. Thus
\[Q=\begin{pmatrix}1/\sqrt2&1/\sqrt6&-1/\sqrt3\\1/\sqrt2&-1/\sqrt6&1/\sqrt3\\0&2/\sqrt6&1/\sqrt3\end{pmatrix},\qquad
R=\begin{pmatrix}\sqrt2&1/\sqrt2&1/\sqrt2\\0&\sqrt{3/2}&1/\sqrt6\\0&0&2/\sqrt3\end{pmatrix}.
\]
Their product recovers $A$, and the diagonal product is two, equal to $|\det A|$ since $\det A=-2$. More generally $|\det Q|=1$ gives the same determinant identity for square $QR$.
\end{solution}
\begin{solution}{linear:ex:projmatrix}
\idea An orthogonal projection fixes its image and makes every residual perpendicular to it.
\writeup The least-squares formula sends $b$ to $A(A^{\mathsf T}A)^{-1}A^{\mathsf T}b$, its closest column-space vector. Conversely if $P^2=P=P^{\mathsf T}$, it fixes $W=\operatorname{Im}P$ and $\langle b-Pb,Pv\rangle=\langle P(b-Pb),v\rangle=0$. This characterises $P_W$. In a basis adapted to $W\oplus W^\perp$, its matrix has $\dim W$ diagonal ones and otherwise zeros, giving trace equal to rank.
\end{solution}
\begin{solution}{linear:ex:parabola}
\idea Discrete orthogonal polynomials diagonalise the fitting equations without changing the model space.
\writeup The data sums are $\sum1=5$, $\sum x^2=10$, $\sum x^4=34$, $\sum y=9$, $\sum xy=2$, $\sum x^2y=30$, with odd pure-power sums zero. The normal equations give $5a+10c=9$, $10b=2$, $10a+34c=30$, hence $(a,b,c)=(3/35,1/5,6/7)$.

The value vectors for $1,x,x^2-2$ are mutually orthogonal, with squared norms $5,10,14$. Their pairings with $y$ are $9,2,12$. Projection therefore gives $9/5+x/5+(6/7)(x^2-2)$, the same quadratic. The diagonal Gram matrix explains why this basis makes the computation simpler.
\end{solution}
\begin{solution}{linear:ex:distpoly}
\idea Set the two residual moments to zero.
\writeup Orthogonality gives $1/3-a-b/2=0$ and $1/4-a/2-b/3=0$, yielding $a=-1/6,b=1$. The residual is $t^2-t+1/6$. With $s=t-1/2$, its squared integral is
\[\int_{-1/2}^{1/2}(s^2-1/12)^2\dd s=1/80-1/144=1/180.
\]
This equals $(\binom42\sqrt5)^{-2}$ from the Hilbert-distance formula.
\end{solution}
\begin{solution}{linear:ex:hilbertint}
\idea Each reciprocal determinant ratio is already an integer binomial expression.
\writeup The ratio equals $(n!)^4/((2n)!^2(2n+1))=((2n+1)\binom{2n}n^2)^{-1}$. Starting with determinant one gives $1/\det H_n=\prod_{k=1}^{n-1}(2k+1)\binom{2k}k^2$. The first four values are $1,12,2160,6048000$, each obtained by this exact product.
\end{solution}
\begin{solution}{linear:ex:obtuse}
\idea A centred simplex constructs the family; perpendicular projection preserves strict obtuseness for the upper bound.
\writeup In $\R^{n+1}$ let $v_i=e_i-(n+1)^{-1}\sum_je_j$. They lie in the dimension-$n$ zero-sum hyperplane and have off-diagonal pairing $-1/(n+1)$. For the upper bound, the case $n=1$ follows because three nonzero real numbers cannot have all pairwise negative products. In dimension $n>1$, suppose there were $n+2$ such vectors and project the first $n+1$ onto the last vector's perpendicular. Their pairings become
\[\langle v_i,v_j\rangle-\frac{\langle v_i,v_m\rangle\langle v_j,v_m\rangle}{\|v_m\|^2}<0.
\]
The projected family is nonzero and has $n+1$ members in dimension $n-1$, contradicting induction. Smaller-family zero-projection edge cases are irrelevant to this contradiction because there are at least three remaining vectors. At $n=0$, one zero vector is allowed by the vacuous pair condition, but two are impossible.
\end{solution}
\begin{solution}{linear:ex:basel}
\idea The sine coefficients of $t$ are multiples of $1/k$, so Bessel bounds their squared sum.
\writeup Product-to-sum integration gives $\int_{-\pi}^\pi\sin(kt)\sin(lt)\dd t=\pi\delta_{kl}$. The normalised sines are orthonormal. Integration by parts gives their coefficient on $t$ as $2\sqrt\pi(-1)^{k+1}/k$. Bessel yields $4\pi\sum_{k\ge1}k^{-2}\le\int_{-\pi}^\pi t^2\dd t=2\pi^3/3$, proving the bound.
\end{solution}
\begin{solution}{linear:ex:legendrenorm}
\idea Endpoint vanishing simplifies Leibniz, and orthogonality reduces the norm to one leading-monomial integral.
\writeup At one, the $n$th derivative of $(t-1)^n(t+1)^n$ has only the Leibniz term $n!2^n$, proving $P_n(1)=1$. Differentiating the leading $t^{2n}$ gives leading coefficient $(2n)!/(2^n(n!)^2)$.

Let $I_n=\int_{-1}^1(1-t^2)^n\dd t$. Integrating the derivative of $t(1-t^2)^n$ gives $(2n+1)I_n=2nI_{n-1}$ for $n\ge1$, with $I_0=2$. This proves $I_n=2^{2n+1}(n!)^2/(2n+1)!$. Integrating Rodrigues's formula against $t^n$ by parts $n$ times yields $\int t^nP_n=I_n/2^n$. Multiplying by the leading coefficient, since all lower degrees are orthogonal to $P_n$, gives $\int P_n^2=2/(2n+1)$.
\end{solution}
\begin{solution}{linear:ex:chebyshev}
\idea Cosine coordinates absorb the weight and turn the rule into a finite trigonometric sum.
\writeup The recurrence $T_{n+1}=2tT_n-T_{n-1}$, starting with $1,t$, gives polynomial degree $n$ and leading coefficient $2^{n-1}$ for $n\ge1$. Substitution $t=\cos\theta$ makes weighted pairings $\int_0^\pi\cos(m\theta)\cos(n\theta)\dd\theta$, proving orthogonality. The zeros have $\theta_k=(2k-1)\pi/(2n)$.

For $0<j<2n$, the geometric sum gives
\[\sum_{k=1}^n\cos(j\theta_k)
 =\Re\frac{1-(-1)^j}{e^{-ij\pi/(2n)}-e^{ij\pi/(2n)}}=0.
\]
The denominator is nonzero and purely imaginary, while the numerator is real. Equal weights $\pi/n$ therefore integrate $T_0$ correctly and every $T_j$ through $2n-1$ to zero. These form a polynomial basis, so the rule has full Gauss exactness; uniqueness of interpolatory weights identifies them as the Gauss weights.
\end{solution}
\begin{solution}{linear:ex:gaussconv}
\idea Exactness removes a polynomial approximant, and positivity bounds both remaining errors by its uniform error.
\writeup The total weight is $M=\int_Iw<\infty$. For $f$ continuous on $[a,b]$, choose a polynomial $q$ with $\|f-q\|_\infty<\varepsilon$ by Weierstrass. Once $2n-1\ge\deg q$, exactness and positive weights give
\[|G_n(f)-\int_Ifw|\le\sum_i\lambda_i|f-q|(t_i)+\int_I|f-q|w\le2M\varepsilon.
\]
Let $n\to\infty$ and then $\varepsilon\downarrow0$. Continuity on the compact closure is what guarantees the uniform approximant used in this proof.
\end{solution}
\begin{solution}{linear:ex:householder}
\idea A reflection reverses its normal and fixes its perpendicular hyperplane; one such move fixes the first basis direction.
\writeup Writing $P_u=uu^{\mathsf T}/\|u\|^2$, we have $P_u^2=P_u=P_u^{\mathsf T}$. Thus $H_u=I-2P_u$ is symmetric and squares to identity. If $u=x-y$ for distinct equal-norm vectors, $\|u\|^2=2\langle x,u\rangle$, giving $H_ux=x-u=y$.

For orthogonal $Q$, if $Qe_1\ne e_1$ use that reflection to obtain $HQe_1=e_1$; otherwise use no reflection. The resulting orthogonal map preserves $e_1^\perp$. Induct there, using at most $n-1$ reflections, each extended by fixing $e_1$. Together with the initial optional one this gives at most $n$, starting from $\pm1$ in dimension one and the empty product in dimension zero. This is the Euclidean Cartan--Dieudonn\'e argument underlying Householder $QR$.
\end{solution}
\begin{solution}{linear:ex:muntz}
\idea Cauchy's product gives the exact Gram ratio for adjoining one power function.
\writeup Power pairings are $\langle t^a,t^b\rangle=1/(a+b+1)$. Their Gram determinant is Cauchy's determinant and is positive for distinct exponents. Adding exponent $m$ multiplies it by $\prod_k(m-a_k)^2/((2m+1)\prod_k(a_k+m+1)^2)$, proving the distance formula.

For sufficiently large $k$, the ratio is $1-(2m+1)/(a_k+m+1)$ in $(0,1)$. The inequality $\log(1-u)\le-u$ gives
\[\log d_N^2\le C-2\sum_{k\text{ sufficiently large}}^N\frac{2m+1}{a_k+m+1}\longrightarrow-\infty,
\]
since $(a_k+m+1)^{-1}\ge(2a_k)^{-1}$ once $a_k\ge m+1$. Thus $d_N\to0$. Monomials whose exponent already occurs need no approximation, and all other monomials satisfy this argument. Polynomial density in $L^2[0,1]$ then gives density of the power span. The converse is a further theorem of M\"untz and Sz\'asz, quoted here.
\end{solution}
\begin{solution}{ex:v8cDa}
\idea Positive row scaling converts weighted fitting to ordinary least squares.
\writeup Minimise $\|W^{1/2}(Ax-y)\|^2$. The equations are $\begin{pmatrix}4&4\\4&6\end{pmatrix}(a,b)^{\mathsf T}=(3,4)^{\mathsf T}$, giving $a=1/4,b=1/2$. The residuals are $(-1/4,1/4,-1/4)$ with weighted sum zero and weighted squared error $1/16+2/16+1/16=1/4$. The positive-definite normal matrix proves uniqueness.
\end{solution}
\begin{solution}{ex:v8cDb}
\idea The feasible affine line has one tangent direction, and the minimum is perpendicular to it.
\writeup Feasibility gives $p=1+ct(1-t)$. Its squared norm is $1+c/3+c^2/30=1/6+(c+5)^2/30$. Thus $c=-5$ uniquely, giving $p=1-5t+5t^2$ and minimum $1/6$. The condition $\langle p,t(1-t)\rangle=0$ is exactly $1/6+c/30=0$, the same minimising equation.
\end{solution}
\begin{solution}{ex:v8cDc}
\idea Symmetry and two moments determine the weights; the squared monic cubic measures the first missed degree.
\writeup The monic cubic is $p_3=t^3-3t/5$, so its nodes are the stated three roots. Symmetry and the moments of one and $t^2$ give $2w+v=2$ and $2w(3/5)=2/3$, hence $w=5/9,v=8/9$. Gauss quadrature proves degree-five exactness. Since $t^6-p_3^2$ has degree four, its error vanishes, while the rule gives zero on $p_3^2$. Thus the $t^6$ error is $\int_{-1}^1p_3^2=8/175$, by Legendre's norm formula or direct integration.
\end{solution}
\begin{solution}{ex:v8cDd}
\idea Self-adjointness of a projection forces the order of the two factors to agree.
\writeup If $PQ$ is an orthogonal projection, $PQ=(PQ)^*=QP$. Conversely commutation gives $(PQ)^2=PQ$ and $(PQ)^*=PQ$, so it is an orthogonal projection. Its image lies in both images, and any vector in their intersection is fixed by the product, proving equality. In $\R^2$, take $P=\begin{pmatrix}1&0\\0&0\end{pmatrix}$ and $Q=\tfrac12\begin{pmatrix}1&1\\1&1\end{pmatrix}$. Their product is neither symmetric nor idempotent: $(PQ)^2=PQ/2\ne PQ$.
\end{solution}

\section*{Sources and further reading}

\begin{itemize}
\item A.-M.~Legendre, \emph{Nouvelles m\'ethodes pour la d\'etermination des orbites des com\`etes}, Paris, 1805, appendix
``Sur la m\'ethode des moindres quarr\'es''.
\item C.~F.~Gauss, \emph{Theoria motus corporum coelestium}, 1809, and \emph{Theoria combinationis observationum erroribus
minimis obnoxiae}, 1821--1823 (the Gauss--Markov theorem); \emph{Methodus nova integralium valores per approximationem
inveniendi}, 1814 (Gauss quadrature).
\item S.~M.~Stigler, \emph{Gauss and the invention of least squares}, Annals of Statistics 9 (1981), 465--474; D.~Teets and
K.~Whitehead, \emph{The discovery of Ceres: how Gauss became famous}, Mathematics Magazine 72 (1999), 83--93.
\item A.-L.~Cauchy, \emph{M\'emoire sur les fonctions altern\'ees et sur les sommes altern\'ees}, Exercices d'analyse et de
physique math\'ematique 2 (1841), 151--159.
\item D.~Hilbert, \emph{Ein Beitrag zur Theorie des Legendre'schen Polynoms}, Acta Mathematica 18 (1894), 155--159;
M.-D.~Choi, \emph{Tricks or treats with the Hilbert matrix}, American Mathematical Monthly 90 (1983), 301--312.
\item P.~Borwein and T.~Erd\'elyi, \emph{The integer Chebyshev problem}, Mathematics of Computation 65 (1996), 661--681.
\item G.~Szeg\H{o}, \emph{Orthogonal Polynomials}, AMS Colloquium Publications 23, 1939 (4th ed.\ 1975), Chapters 2--3.
\item G.~H.~Golub and J.~H.~Welsch, \emph{Calculation of Gauss quadrature rules}, Mathematics of Computation 23 (1969),
221--230: Proposition~\ref{linear:prop:jacobi} as an algorithm.
\item P.~Jordan and J.~von Neumann, \emph{On inner products in linear, metric spaces}, Annals of Mathematics 36 (1935),
719--723.
\item C.~H.~M\"untz, \emph{\"Uber den Approximationssatz von Weierstrass}, Schwarz-Festschrift, Berlin, 1914, 303--312.
\item L.~N.~Trefethen and D.~Bau, \emph{Numerical Linear Algebra}, SIAM, 1997, Lectures 7--11 (QR, Householder, least
squares) and Lecture~19 (stability of least squares).
\item S.~Axler, \emph{Linear Algebra Done Right}, 4th ed., Springer, 2024, Chapters 6--7.
\end{itemize}
